% Source: OpenAI2026AreaLaw, 09-amplification.tex, lines 203--215,
% openai/math@adc7f1241b42e322a6451854ab7e4b4c146bf78a.
% The physical module is ChannelWordOmission at
% d53502c3def2a0b8abd0d1685868ffcc9eef0d13.
% This leaf proves the omitted-prefix and expected time-integral steps.
% Spatial one-event estimates and summation of their distance tail are separate.

\section{Omitting events from physical channel words}
\label{sec:al_channel_word_omission}

Let $I,J$ be finite sets, let $q\in\N$, and put
$\mathcal H=(\C^q)^{\otimes I}$ and $\mathcal K=\C^J$.
The spectator space $\mathcal K$ may have dimension zero.
For each label $i$ in an event alphabet $F$, let $k_i$ be a
physical matrix on $\mathcal H$. For a matrix $A$, the notation
$A^{1/2}$ means its positive square root when $A\ge0$, and zero
otherwise. Write
\begin{align}
    \label{eq:al_omission_root_map}
    \mathcal E_i(B)
     & =\bigl((\mathbb 1-k_i)^{1/2}\otimes\mathbb 1\bigr)
    B\bigl((\mathbb 1-k_i)^{1/2}\otimes\mathbb 1\bigr)
    +\bigl(k_i^{1/2}\otimes\mathbb 1\bigr)
    B\bigl(k_i^{1/2}\otimes\mathbb 1\bigr).
\end{align}
Thus $\mathcal E_i$ is the actual spectator root channel from
Definition~\ref{def:al_local_root_channel}. For a chronological word
$w=(i_0,\ldots,i_{m-1})$, set
$\mathcal E_w=\mathcal E_{i_{m-1}}\circ\cdots\circ\mathcal E_{i_0}$,
with $\mathcal E_{\varnothing}=\mathrm{id}$.
All norms below are Euclidean operator norms, and every input
$B$ is an arbitrary complex matrix on $\mathcal H\otimes\mathcal K$.

\begin{theorem}[Subtraction and contraction along a word]
    \label{thm:al_channel_word_difference_contraction}
    \uses{def:al_channel_word}
    \lean{TNLean.PEPS.AreaLaw.spectatorRootChannelWord_sub,
        TNLean.PEPS.AreaLaw.norm_spectatorRootChannelWord_sub_le}
    \leanok
    For all words $w$ and operators $B,C$,
    \begin{align}
        \label{eq:al_omission_word_subtraction}
        \mathcal E_w(B-C) & =\mathcal E_w(B)-\mathcal E_w(C).
    \end{align}
    If $0\le k_i\le\mathbb 1$ for all $i\in F$, then
    \begin{align}
        \label{eq:al_omission_difference_contraction}
        \|\mathcal E_w(B)-\mathcal E_w(C)\| & \le\|B-C\|.
    \end{align}
\end{theorem}
\begin{proof}
    \leanok
    \uses{def:al_local_root_channel,thm:al_channel_word_norm}
    Both summands in~(\ref{eq:al_omission_root_map}) preserve
    subtraction. Induction through the chronological composition
    gives~(\ref{eq:al_omission_word_subtraction}). Apply
    Theorem~\ref{thm:al_channel_word_norm} to $B-C$ to obtain
    (\ref{eq:al_omission_difference_contraction}).
\end{proof}

Let $P$ be a predicate on $F$. Denote by $w|_P$ the subword
consisting of the labels satisfying $P$, in the same order and
with the same multiplicities. The full and retained evolutions
are evaluated from the same original word $w$.

\begin{definition}[Actual omission defect]
    \label{def:al_channel_word_omission_defect}
    \uses{def:al_channel_word}
    \lean{TNLean.PEPS.AreaLaw.channelWordOmissionDefect}
    \leanok
    For an initial operator $B$, define
    \begin{align}
        \label{eq:al_omission_defect}
        D_P(w;B) & =\|\mathcal E_w(B)-\mathcal E_{w|_P}(B)\|.
    \end{align}
    In particular $D_P(\varnothing;B)=0$.
\end{definition}

\begin{theorem}[Uniform bound for the actual defect]
    \label{thm:al_channel_word_omission_bounds}
    \uses{def:al_channel_word_omission_defect}
    \lean{TNLean.PEPS.AreaLaw.channelWordOmissionDefect_bounds}
    \leanok
    Suppose $0\le k_i\le\mathbb 1$ for all $i\in F$.
    For every word $w$ and operator $B$,
    \begin{align}
        \label{eq:al_omission_uniform_bound}
        0\le D_P(w;B) & \le 2\|B\|.
    \end{align}
\end{theorem}
\begin{proof}
    \leanok
    \uses{thm:al_channel_word_norm}
    Nonnegativity follows from the norm. The triangle inequality
    and contraction of both chronological words give
    $D_P(w;B)\le\|\mathcal E_w(B)\|+
        \|\mathcal E_{w|_P}(B)\|\le2\|B\|$.
\end{proof}

\begin{theorem}[One-event increment at the retained prefix]
    \label{thm:al_channel_word_omission_increment}
    \uses{def:al_channel_word_omission_defect}
    \lean{TNLean.PEPS.AreaLaw.channelWordOmissionDefect_append_sub_le}
    \leanok
    Suppose $0\le k_i\le\mathbb 1$ for all $i\in F$.
    Write $wi$ for appending the label $i$ to $w$. Then
    \begin{align}
        \label{eq:al_omission_one_increment}
        D_P(wi;B)-D_P(w;B)
         & \le
        \left\{\begin{array}{ll}
                   0,                        & P(i),      \\
                   \|(\mathcal E_i-\mathrm{id})
                   (\mathcal E_{w|_P}(B))\|, & \neg P(i).
               \end{array}\right.
    \end{align}
\end{theorem}
\begin{proof}
    \leanok
    \uses{def:al_local_root_channel,thm:al_root_channel_contraction,
        thm:al_channel_word_append}
    Put $A=\mathcal E_w(B)$ and $R=\mathcal E_{w|_P}(B)$.
    If $P(i)$ holds, both evolutions append $\mathcal E_i$, so
    $D_P(wi;B)=\|\mathcal E_i(A)-\mathcal E_i(R)\|
        \le\|A-R\|=D_P(w;B)$.
    If $\neg P(i)$, the retained evolution stays at $R$. Hence
    \begin{align}
        \label{eq:al_omission_triangle_step}
        D_P(wi;B)
         & =\|\mathcal E_i(A)-R\|\nonumber        \\
         & \le\|\mathcal E_i(A)-\mathcal E_i(R)\|
        +\|\mathcal E_i(R)-R\|\nonumber           \\
         & \le D_P(w;B)+\|\mathcal E_i(R)-R\|.
    \end{align}
    Subtract $D_P(w;B)$ to obtain
    (\ref{eq:al_omission_one_increment}).
\end{proof}

\begin{theorem}[Literal omitted-prefix telescope]
    \label{thm:al_channel_word_omission_telescope}
    \uses{def:al_channel_word_omission_defect}
    \lean{TNLean.PEPS.AreaLaw.channelWordOmissionDefect_le_sum_prefix}
    \leanok
    Suppose $0\le k_i\le\mathbb 1$ for all $i\in F$.
    For a word $w=(i_0,\ldots,i_{m-1})$, let
    $w_{<j}=(i_0,\ldots,i_{j-1})$ be its actual earlier prefix.
    Then
    \begin{align}
        \label{eq:al_omission_prefix_telescope}
        D_P(w;B) & \le
        \sum_{0\le j<m,\ \neg P(i_j)}
        \|(\mathcal E_{i_j}-\mathrm{id})
        (\mathcal E_{w_{<j}|_P}(B))\|.
    \end{align}
    Repeated labels contribute at each of their omitted positions.
\end{theorem}
\begin{proof}
    \leanok
    \uses{thm:al_channel_word_omission_increment}
    The real-valued prefix increments telescope:
    \begin{align}
        \label{eq:al_omission_real_telescope}
        D_P(w;B) & =
        \sum_{j=0}^{m-1}
        \bigl(D_P(w_{<j}i_j;B)-D_P(w_{<j};B)\bigr),
    \end{align}
    since the initial defect is zero. Apply
    (\ref{eq:al_omission_one_increment}) at each actual prefix
    $w_{<j}$. A retained position contributes at most zero; an
    omitted position contributes its action on the retained
    earlier prefix. Summing proves
    (\ref{eq:al_omission_prefix_telescope}).
\end{proof}

Now assume $F$ is finite, set $F_P=\{i\in F:P(i)\}$, and write
$\mathcal E_u^P$ for the chronological composition of the
restricted effects for a word $u$ over $F_P$. For any finite
alphabet $A$ and $t\ge0$, let $\nu_{A,t}$ be the probability
measure assigning each individual word $u$ of length $m$ the mass
\begin{align}
    \label{eq:al_omission_word_mass}
    \nu_{A,t}(\{u\}) & =e^{-|A|t}\frac{t^m}{m!}.
\end{align}
At $t=0$, or when $A$ is empty, this measure is concentrated on
the empty word.

\begin{theorem}[Expected omission error from retained one-event errors]
    \label{thm:al_channel_word_expected_omission}
    \uses{def:al_channel_word_omission_defect}
    \lean{TNLean.PEPS.AreaLaw.integrable_and_integral_channelWordOmissionDefect_le}
    \leanok
    Suppose $0\le k_i\le\mathbb 1$ for all $i\in F$.
    For every $T\ge0$ and every initial operator $B$, the actual
    omission defect $w\mapsto D_P(w;B)$ is integrable under
    $\nu_{F,T}$, and
    \begin{align}
        \label{eq:al_omission_expected_time_integral}
        \int D_P(w;B)\,d\nu_{F,T}(w)
         & \le\int_{(0,T]}\sum_{i\in F,\ \neg P(i)}
        \left(\int
        \|(\mathcal E_i-\mathrm{id})(\mathcal E_u^P(B))\|
        \,d\nu_{F_P,s}(u)\right)\,ds.
    \end{align}
    The left side is the expectation of the operator norm of the
    difference for the full word and its own retained subword.
\end{theorem}
\begin{proof}
    \leanok
    \uses{thm:al_channel_word_omission_bounds,
        thm:al_channel_word_omission_telescope,thm:al_channel_word_filter,
        thm:al_channel_word_norm,thm:al_root_channel_contraction}
    Define the nonnegative retained-word cost
    \begin{align}
        \label{eq:al_omission_retained_cost}
        c_i(u) & =\|(\mathcal E_i-\mathrm{id})(\mathcal E_u^P(B))\|.
    \end{align}
    The triangle inequality and contraction imply
    $0\le c_i(u)\le2\|\mathcal E_u^P(B)\|\le2\|B\|$.
    The defect is also bounded by $2\|B\|$ by
    (\ref{eq:al_omission_uniform_bound}), and so it is integrable
    under the probability measure on the countable word space.
    The omitted-prefix sum in
    (\ref{eq:al_omission_prefix_telescope}) is bounded by
    $2\|B\|\,|w|$ and is integrable, since
    $\int|w|\,d\nu_{F,T}(w)=|F|T$.

    Let $\pi_P(w_{<j})$ be the actual retained prefix, now viewed
    as a word over $F_P$. Its observable is
    $\mathcal E_{w_{<j}|_P}(B)=\mathcal E^P_{\pi_P(w_{<j})}(B)$.
    The retained-prefix occupation identity for the atomic
    word law gives
    \begin{align}
        \label{eq:al_omission_occupation_identity}
         & \int\sum_{0\le j<|w|,\ \neg P(i_j)}
        c_{i_j}(\pi_P(w_{<j}))\,d\nu_{F,T}(w)\nonumber  \\
         & \qquad=\int_{(0,T]}\sum_{i\in F,\ \neg P(i)}
        \left(\int c_i(u)\,d\nu_{F_P,s}(u)\right)\,ds.
    \end{align}
    Integrate~(\ref{eq:al_omission_prefix_telescope}) and apply
    (\ref{eq:al_omission_occupation_identity}). This proves
    (\ref{eq:al_omission_expected_time_integral}) using the actual
    retained-prefix law.
\end{proof}

Theorem~\ref{thm:al_channel_word_expected_omission} is the
time-integration step in \cite[\texttt{09-amplification.tex},
    lines 203--215]{OpenAI2026AreaLaw}. The spatial bound for each
omitted event and the subsequent sum over anchor distances are
separate estimates. The result here applies for arbitrary
finite spectators and arbitrary complex $B$, and requires no
assumption that a sampled prefix is independent of its word.
